\subsection{Weaknesses and Limitations}
\label{subsec:schwaechen-grenzen}

Technical weaknesses, deliberate scope boundaries, and unverified areas must
be distinguished. Only the first category directly denotes a disadvantage of
the implementation.
Although scope boundaries create a need for adaptation, they do not constitute
an unmet requirement. A lack of verification does not yet permit either a
positive or a negative technical judgment.

Consolidation in the master repository increases internal complexity: changes
to the core, generator, or catalog require regression tests across several
profiles. In addition, the frontend and backend do not use the shared JSON
schemas at runtime; contract changes remain manual. Direct CLI selection
distinguishes \texttt{optional} and \texttt{selectable} less clearly than the
dialog does. Both discrepancies create risks of maintenance errors or
incorrect operation.
The examined HEAD also has conflicting evidence: the web profiles and the local
tooling run passed, while the external tooling, desktop-profile, and Windows
checks failed.
This finding does not retroactively call the earlier, commit-specific
acceptance into question, but it prevents a fully successful assessment of the
more recent state.

The deliberate exclusions are a backend in \texttt{web-only}, a cloud API in
\texttt{desktop-local}, and PostgreSQL without a backend. Domain logic,
authentication, role models, mandatory persistence, cloud providers, and
native domain-specific commands likewise remain product-specific. The
new decision framework for product persistence provides no universal
local-file, embedded-database, or synchronization implementation. This
deliberate scope boundary is an extension path, not an implementation defect
\parencite{kleiveist2026persistence}.

Compose startup, production deployment, Release Validation, and Branch
Protection have not been conclusively verified. Unsigned Linux and macOS
packages were built, while Windows packaging failed. Signing, notarization,
publication, and production cloud operation remain product- or
infrastructure-specific
\parencite{kleiveist2026ci-core-run,kleiveist2026ci-desktop-run,kleiveist2026acceptance}.
The packaging results therefore demonstrate neither signed releases nor
production cloud operation. The subject of the evaluation is the template
baseline, not the production readiness of a product.
The remaining risks therefore primarily concern operations, platforms, and
release processes.
\beginnerinput{workspace/statement/700_bewertung/730}
